\subsection{Extra instances do not remove a shared licensing dependency}
\label{exec17:s8-availability}

The inspected Bayesia API description separates model-file location from
license validation: a Java configuration can retain model files on customer
servers while validating credentials with the vendor's license server when
launching an instance. It states a permanent connection requirement. This
has a specific acquisition consequence. A proposal for several local replicas
must explain which dependencies those replicas share. Local data residency
does not, by itself, imply independent recovery or offline operation
\citep{exec13_bayesia_api}.

Consider a deliberately narrow recovery experiment. All replicas are initially
unloaded, and a required model must become ready during a stated recovery
window. Suppose validation must be available in that window, and at least one
local worker must be able to load the model. Let $C$ indicate common validation
availability and $R_j$ local readiness for replica $j$. Write
$\Prob(C=0)=q$ and $\Prob(R_j=0)=p$. In this authored model the $R_j$ are
mutually independent and independent of $C$. The success event is
\begin{equation}
 \{C=1\}\cap\bigcup_{j=1}^m\{R_j=1\},\qquad
 A_m=(1-q)(1-p^m).
 \label{exec17:s8-common-recovery}
\end{equation}
The formula follows by taking the complement of all local workers failing,
then applying the stated independence from validation. It concerns a cold
recovery request. It is not a statement that disconnecting the license server
terminates a currently loaded engine or makes every ordinary query fail;
that behavior is not established by the inspected page.

With stipulated $p=1/10$ and $q=1/50$, one, two and three replicas have success
probabilities .882, .9702 and .97902. Ignoring the shared validation gate would
report .9, .99 and .999 instead. The difference matters to a capacity purchase:
a two-replica calculation that drops the common dependency appears to meet a
.99 target but actually does not meet it in this very model. Increasing $m$
cannot solve that problem, since
\begin{equation}
 A_m\leq1-q=.98<.99.
 \label{exec17:s8-recovery-ceiling}
\end{equation}
The probability calculation is reproduced by summing every common/local
readiness combination for each of the three capacities
\citep{exec17_s8_availability}.

For a general desired probability $a$, the same model requires
$q\leq1-a$ before replica count can be sufficient. When $q<1-a$ and
$0<p<1$, a finite sufficient count satisfies
\begin{equation}
 p^m\leq1-\frac{a}{1-q}.
 \label{exec17:s8-recovery-condition}
\end{equation}
At $q=1-a$, the right side is zero, so no finite count with $p>0$ attains
the boundary. These conditions do not estimate $p$ or $q$; they identify which
measurement or architectural change could alter a purchase decision. A second
local machine cannot resolve an unavailable common credential path. A different
permitted validation arrangement, a genuinely ready standby or an alternative
engine may change the event structure, but each requires its own documented
rights, scientific equivalence and operating evidence.

The standby distinction is especially easy to obscure in a benchmark.
Keeping a model already loaded changes the recovery experiment and can consume
a licensed instance during the quiet period. Starting the standby only after
failure changes its readiness time and can expose the shared validation path.
Neither arrangement is universally cheaper or more reliable from the public
description alone. A quote should specify whether the required capacity
includes retained standby models, when validation occurs and what recovery
behavior the vendor actually supports. A performance comparison should count
the preparation and occupancy of that exact arrangement.

The independence assumption also needs scrutiny. Replicas on one host can
share a process supervisor, memory pressure and storage. Replicas on different
hosts can still share a network route, model store or license service. If local
readiness is perfectly common rather than independent, adding replicas leaves
local success at $1-p$, making the cold-recovery probability $(1-q)(1-p)$
for every $m$. The appropriate input is a joint failure model, not a list of
machines. This is a practical reason to preserve configuration and dependency
identity alongside observed failure counts.

The inexpensive prototype can first exercise the common layer against a
declared unavailable-validation response and an interrupted model load. It
should retain the recovery request, each attempt, elapsed time and the eventual
fallback or unresolved result. Such fixtures can demonstrate correct accounting
and refusal before a paid engine is available. They cannot establish actual
license-server uptime or vendor failover behavior. A later permitted trial
should inspect those behaviors using the exact contracted edition and recovery
arrangement. A broad availability claim waits for observations from the relevant
workload and time horizon.

This dependency calculation belongs in the commercial comparison because it
can change the item being purchased. A capacity-only quotation may offer more
instances without addressing the binding common recovery condition. A proposed
open or commercial alternative must preserve the same causal outputs before
its different recovery architecture is economically comparable. The useful
procurement conclusion is conditional and specific: price the supported
configuration that can satisfy the required recovery experiment, and record
when the current evidence cannot identify such a configuration.
